% unit: u02
% sources: sources/068-an-overview-of-key-ideas.txt, sources/120-an-overview-of-linear-algebra.txt, sources/121-mit18-06scf11-ses1-13sum-pdf-download-file.txt, sources/122-an-overview-of-key-ideas.txt, short_plan.json (s01)
\section{An overview of key ideas}\label{sec:u02}

Before the course settles into elimination, it is worth seeing the whole road at once.
Linear algebra moves from \emph{vectors} to \emph{matrices} to \emph{subspaces}: what you do
with vectors is take combinations; a matrix packages those combinations as a single
operation $\vb{x}\mapsto A\vb{x}$; and the set of all combinations is a subspace. This
section, adapted from Strang's opening review in 18.085, follows two $3\times3$ matrices that
differ in a single column. One is invertible and its inverse can be written down by hand; the
other has dependent columns, solves $C\vb{x}=\vb{b}$ only for special $\vb{b}$, and shows
what a subspace is. Almost every idea of Part~I appears here in preview, and later sections
make each one precise.

\subsection{Vectors: take their combinations}

Work in $\R^3$, the space of vectors with three real components, with
\[
  \vb{u}=\begin{bmatrix} 1\\ -1\\ 0\end{bmatrix},\qquad
  \vb{v}=\begin{bmatrix} 0\\ 1\\ -1\end{bmatrix},\qquad
  \vb{w}=\begin{bmatrix} 0\\ 0\\ 1\end{bmatrix}.
\]
The only operations we need are multiplying a vector by a scalar and adding vectors, which
together produce linear combinations $x_1\vb{u}+x_2\vb{v}+x_3\vb{w}$
(Definition~\ref{def:u01-combination}). All multiples $c\,\vb{u}$ fill a line through the
origin. All multiples of $\vb{v}$ fill a second line. Because $\vb{v}$ is not a multiple of
$\vb{u}$, the combinations $x_1\vb{u}+x_2\vb{v}$ fill something bigger than either line: a
plane through the origin. Whether adding $\vb{w}$ fills all of $\R^3$ is the question the
rest of the section answers.

\subsection{Matrices: the difference matrix and its inverse}

Put the three vectors into the columns of a matrix. Then the combination with weights
$x_1,x_2,x_3$ is exactly the matrix-vector product of Definition~\ref{def:u01-Ax}:
\begin{equation}\label{eq:u02-A}
  A\vb{x}=\begin{bmatrix} 1 & 0 & 0\\ -1 & 1 & 0\\ 0 & -1 & 1\end{bmatrix}
  \begin{bmatrix} x_1\\ x_2\\ x_3\end{bmatrix}
  = x_1\vb{u}+x_2\vb{v}+x_3\vb{w}
  = \begin{bmatrix} x_1\\ x_2-x_1\\ x_3-x_2\end{bmatrix}.
\end{equation}
The calculation is the same as before but the point of view has changed. Before, numbers
$x_j$ multiplied vectors; now the matrix $A$ acts on the vector $\vb{x}$ and produces an
output. This $A$ is a \emph{difference matrix}: each output component is a difference of
consecutive inputs (with $x_0=0$).

\begin{example}[Differences of squares]\label{ex:u02-squares}
Feed in the squares $\vb{x}=(1,4,9)$:
\[
  A\begin{bmatrix} 1\\ 4\\ 9\end{bmatrix}
  =\begin{bmatrix} 1\\ 4-1\\ 9-4\end{bmatrix}=\begin{bmatrix} 1\\ 3\\ 5\end{bmatrix}.
\]
Differences of consecutive squares are the odd numbers, since $(k+1)^2-k^2=2k+1$.
\end{example}

The deeper question runs backwards: given the output $\vb{b}$, find the input $\vb{x}$ with
$A\vb{x}=\vb{b}$. The three equations are
\[
  x_1=b_1,\qquad x_2-x_1=b_2,\qquad x_3-x_2=b_3 .
\]
Because $A$ is lower triangular they can be solved from the top down: $x_1=b_1$, then
$x_2=b_1+b_2$, then $x_3=b_1+b_2+b_3$. Read the answer as a matrix times $\vb{b}$: the
coefficients of $b_1$ form the first column, and so on.
\begin{equation}\label{eq:u02-S}
  \vb{x}=\begin{bmatrix} b_1\\ b_1+b_2\\ b_1+b_2+b_3\end{bmatrix}
  =\begin{bmatrix} 1 & 0 & 0\\ 1 & 1 & 0\\ 1 & 1 & 1\end{bmatrix}
  \begin{bmatrix} b_1\\ b_2\\ b_3\end{bmatrix} = S\vb{b}.
\end{equation}
The matrix $S$ takes running sums. It undoes $A$, and $A$ undoes $S$:

\begin{proposition}[The sum matrix inverts the difference matrix]\label{prop:u02-sum}
For every $\vb{x}\in\R^3$, $S(A\vb{x})=\vb{x}$, and for every $\vb{b}\in\R^3$,
$A(S\vb{b})=\vb{b}$. Consequently $A\vb{x}=\vb{b}$ has exactly one solution for every
$\vb{b}$, namely $\vb{x}=S\vb{b}$, and $A\vb{x}=\vb{0}$ only for $\vb{x}=\vb{0}$.
\end{proposition}

\begin{proof}
Running sums of differences telescope:
\[
  S(A\vb{x}) = S\begin{bmatrix} x_1\\ x_2-x_1\\ x_3-x_2\end{bmatrix}
  = \begin{bmatrix} x_1\\ x_1+(x_2-x_1)\\ x_1+(x_2-x_1)+(x_3-x_2)\end{bmatrix}
  = \begin{bmatrix} x_1\\ x_2\\ x_3\end{bmatrix}.
\]
Differences of running sums return the original entries:
$A(S\vb{b}) = \bigl(b_1,\ (b_1+b_2)-b_1,\ (b_1+b_2+b_3)-(b_1+b_2)\bigr)=(b_1,b_2,b_3)$.
The second identity shows that $\vb{x}=S\vb{b}$ solves $A\vb{x}=\vb{b}$. If $\vb{x}$ is any
solution, then $\vb{x}=S(A\vb{x})=S\vb{b}$ by the first identity, so the solution is unique.
With $\vb{b}=\vb{0}$ this unique solution is $S\vb{0}=\vb{0}$.
\end{proof}

A matrix that undoes $A$ in this way is called the \emph{inverse} of $A$ and written
$A^{-1}$; here $A^{-1}=S$. Section~\ref{sec:u04} defines inverses in general and shows how to
compute them. Think of $A$ as a transform taking $\vb{x}$ to $\vb{b}$ and of $S$ as the
inverse transform taking $\vb{b}$ back to $\vb{x}$. For example $S(1,3,5)=(1,1+3,1+3+5)=(1,4,9)$
recovers the squares. Calculus has the same structure: differences are the discrete version of
the derivative, running sums the discrete version of the integral, and
Proposition~\ref{prop:u02-sum} is a small fundamental theorem of calculus.

\subsection{The cyclic difference matrix}

Keep $\vb{u}$ and $\vb{v}$ but replace $\vb{w}$ by $\vb{w}^{*}=(-1,0,1)$. The new matrix
\begin{equation}\label{eq:u02-C}
  C\vb{x}=\begin{bmatrix} 1 & 0 & -1\\ -1 & 1 & 0\\ 0 & -1 & 1\end{bmatrix}
  \begin{bmatrix} x_1\\ x_2\\ x_3\end{bmatrix}
  =\begin{bmatrix} x_1-x_3\\ x_2-x_1\\ x_3-x_2\end{bmatrix}
\end{equation}
still takes differences, but now cyclically: the first row compares $x_1$ with $x_3$, as if
the three entries sat on a circle. Only one column changed, yet $C$ behaves completely
differently from $A$.

\emph{Zero output from a nonzero input.} If $x_1=x_2=x_3=c$, every difference is zero, so
$C(c,c,c)=\vb{0}$. A whole line of inputs, all multiples of $(1,1,1)$, is sent to zero.

\begin{proposition}[No inverse when a nonzero vector goes to zero]\label{prop:u02-noinverse}
If $C\vb{x}=\vb{0}$ for some $\vb{x}\neq\vb{0}$, then there is no matrix $T$ with
$T(C\vb{y})=\vb{y}$ for all $\vb{y}$; that is, $C$ has no inverse.
\end{proposition}

\begin{proof}
Such a $T$ would give $\vb{x}=T(C\vb{x})=T\vb{0}=\vb{0}$, since any matrix times the zero
vector is the zero vector. This contradicts $\vb{x}\neq\vb{0}$. Once $C$ has crushed
$\vb{x}$ to zero, nothing can bring it back.
\end{proof}

\emph{Solvable only for special right-hand sides.} Write out $C\vb{x}=\vb{b}$ and add the
three equations. The left sides cancel, $(x_1-x_3)+(x_2-x_1)+(x_3-x_2)=0$, so the right
sides must satisfy $0=b_1+b_2+b_3$. This condition is not only necessary but sufficient.

\begin{proposition}[When the cyclic system is solvable]\label{prop:u02-solvable}
$C\vb{x}=\vb{b}$ has a solution if and only if $b_1+b_2+b_3=0$. In that case the solutions
are exactly
\[
  \vb{x} = \begin{bmatrix} 0\\ b_2\\ b_2+b_3\end{bmatrix} + c\begin{bmatrix} 1\\ 1\\ 1\end{bmatrix},
  \qquad c\in\R .
\]
\end{proposition}

\begin{proof}
Necessity was shown above by adding the equations. Conversely, suppose $b_1+b_2+b_3=0$ and
let $\vb{x}_p=(0,\,b_2,\,b_2+b_3)$. Then
\[
  C\vb{x}_p = \begin{bmatrix} 0-(b_2+b_3)\\ b_2-0\\ (b_2+b_3)-b_2\end{bmatrix}
  = \begin{bmatrix} b_1\\ b_2\\ b_3\end{bmatrix},
\]
using $-(b_2+b_3)=b_1$. Adding $c(1,1,1)$ does not change $C\vb{x}$, because
$C(c,c,c)=\vb{0}$. Finally, if $\vb{x}$ is any solution, then $\vb{y}=\vb{x}-\vb{x}_p$ satisfies
$C\vb{y}=\vb{b}-\vb{b}=\vb{0}$, so $y_2-y_1=0$ and $y_3-y_2=0$; hence $\vb{y}=c(1,1,1)$ for
some $c$, and every solution has the stated form.
\end{proof}

In a physical model where $x_1,x_2,x_3$ are positions of masses joined in a ring by springs
and $\vb{b}$ holds the applied forces, the condition $b_1+b_2+b_3=0$ says the forces must
balance; the free constant $c$ says the whole ring can be shifted without changing any
spring.

\begin{example}[Solving and failing with $C$]\label{ex:u02-cyclic}
Take $\vb{b}=(1,2,-3)$. Its components sum to $0$, so the system is solvable, and
Proposition~\ref{prop:u02-solvable} gives $\vb{x}_p=(0,\,2,\,2-3)=(0,2,-1)$. Check:
$C\vb{x}_p=(0-(-1),\ 2-0,\ -1-2)=(1,2,-3)$. All solutions are $(0,2,-1)+c(1,1,1)$, for
instance $(1,3,0)$ with $c=1$. For $\vb{b}=(1,1,1)$ the components sum to $3\neq0$, and
$C\vb{x}=\vb{b}$ has no solution at all.
\end{example}

\subsection{Subspaces: what all the combinations fill}

Which vectors $\vb{b}$ are combinations of the columns $\vb{u},\vb{v},\vb{w}^{*}$ of $C$?
Each column has components adding to zero, and that property survives addition and
scalar multiplication, so every combination $C\vb{x}$ lies on the plane
$b_1+b_2+b_3=0$. Proposition~\ref{prop:u02-solvable} says that every vector on that plane is
reached. So \emph{all vectors $C\vb{x}$} form exactly the plane $b_1+b_2+b_3=0$. In fact
$\vb{w}^{*}=-(\vb{u}+\vb{v})$ lies in the plane of $\vb{u}$ and $\vb{v}$ and adds nothing
new (Proposition~\ref{prop:u01-dependent}). For $A$, Proposition~\ref{prop:u02-sum} says
every $\vb{b}$ in $\R^3$ is reached, so the combinations of $\vb{u},\vb{v},\vb{w}$ fill all of
$\R^3$. Figure~\ref{fig:u02-plane} shows the difference.

\begin{figure}[ht]
\centering
\begin{tikzpicture}[>=Stealth]
\begin{axis}[view={44}{16}, width=11cm, height=9cm, axis x line=none, axis y line=none, axis z line=center, axis line style={gray!70},
  ticks=none, zlabel={$b_3$},
  xmin=-1.7, xmax=1.7, ymin=-1.7, ymax=1.7, zmin=-2.7, zmax=2.7, clip=false]
  \addplot3[surf, shader=faceted, colormap={single}{color=(blue!12) color=(blue!12)},
    faceted color=blue!45, opacity=0.65, samples=7, domain=-1.3:1.3, y domain=-1.3:1.3] {-x-y};
  \draw[->,very thick] (0,0,0) -- (1,-1,0) node[below left] {$\vb{u}$};
  \draw[->,very thick] (0,0,0) -- (0,1,-1) node[below right] {$\vb{v}$};
  \draw[->,very thick,blue] (0,0,0) -- (-1,0,1) node[left] {$\vb{w}^{*}$};
  \draw[->,very thick,red!75!black] (0,0,0) -- (0,0,1) node[right] {$\vb{w}$};
  \draw[thick,densely dotted,red!75!black] (0,0,1) -- (-0.333,-0.333,0.667);
  \fill[red!75!black] (-0.333,-0.333,0.667) circle (1.6pt);
  \node[blue!70!black,font=\small] at (1.3,1.3,-2.6) [below] {$b_1+b_2+b_3=0$};
\end{axis}
\end{tikzpicture}
\caption{The columns $\vb{u}=(1,-1,0)$, $\vb{v}=(0,1,-1)$ and $\vb{w}^{*}=(-1,0,1)$ of the
cyclic matrix $C$ all lie in the plane $b_1+b_2+b_3=0$ (shaded), so every $C\vb{x}$ stays in
that plane. The third column $\vb{w}=(0,0,1)$ of the difference matrix $A$ points out of the
plane (the dotted segment drops from its tip to the nearest point of the plane), so the
combinations of $\vb{u},\vb{v},\vb{w}$ fill all of $\R^3$.}\label{fig:u02-plane}
\end{figure}

Here are the words for what we have seen; Sections~\ref{sec:u06}, \ref{sec:u07} and
\ref{sec:u10} develop them in full.

\begin{definition}[Independent and dependent vectors]\label{def:u02-independent}
Vectors $\vb{v}_1,\dots,\vb{v}_n$ are \emph{dependent} if some combination
$c_1\vb{v}_1+\cdots+c_n\vb{v}_n$ with weights not all zero equals the zero vector. Otherwise
they are \emph{independent}: the only combination giving $\vb{0}$ is the one with all
weights zero.
\end{definition}

In matrix language, the columns of a matrix are independent exactly when $A\vb{x}=\vb{0}$
forces $\vb{x}=\vb{0}$. So $\vb{u},\vb{v},\vb{w}$ are independent
(Proposition~\ref{prop:u02-sum}), while $\vb{u}+\vb{v}+\vb{w}^{*}=C(1,1,1)=\vb{0}$ shows that
$\vb{u},\vb{v},\vb{w}^{*}$ are dependent.

\begin{definition}[Vector space, subspace, basis]\label{def:u02-subspace}
A \emph{vector space} is a collection of vectors closed under linear combinations: any
combination of vectors in the collection is again in it. A \emph{subspace} is a vector space
contained in a larger one. A \emph{basis} of $\R^n$ is a set of $n$ independent vectors in
$\R^n$.
\end{definition}

The plane $b_1+b_2+b_3=0$ is a subspace of $\R^3$: it is the set of all combinations of
$\vb{u}$ and $\vb{v}$, and a combination of combinations is a combination. A filled triangle
with corners $\vb{0},\vb{u},\vb{v}$ is \emph{not} a subspace, because $2\vb{u}$ leaves it.
Taking all combinations of some vectors is the standard way to build a subspace, and it
always produces a subspace. The subspaces of $\R^3$ are: the origin alone, the lines
through the origin, the planes through the origin, and $\R^3$ itself, of dimensions
$0,1,2,3$. The zero vector belongs to every one of them, since $0\cdot\vb{v}=\vb{0}$.

For $n$ vectors in $\R^n$ there are three equivalent ways to say ``basis'', which this
section has seen side by side for $\vb{u},\vb{v},\vb{w}$: (i) the vectors are independent;
(ii) their combinations fill $\R^n$; (iii) the $n\times n$ matrix with those columns is
invertible. The vectors $\vb{u},\vb{v},\vb{w}^{*}$ fail all three. The proof that the three
statements always agree comes in Section~\ref{sec:u10}.

\begin{example}[A basis of the plane from Section~\ref{sec:u01}]\label{ex:u02-basis2}
The columns $\vb{a}_1=(2,-1)$ and $\vb{a}_2=(-1,2)$ of the first example of the course form a
basis of $\R^2$. They are independent: if $c_1\vb{a}_1+c_2\vb{a}_2=\vb{0}$ then
$2c_1-c_2=0$ and $-c_1+2c_2=0$; the first gives $c_2=2c_1$, and the second becomes
$-c_1+4c_1=3c_1=0$, so $c_1=c_2=0$. Their combinations fill $\R^2$: we found
$1\,\vb{a}_1+2\,\vb{a}_2=(0,3)$, and $\vb{a}_1+\vb{a}_2=(1,1)$, and
Exercise~\ref{ex:u01-every-b} gives the weights for every $\vb{b}$. And the matrix
$\begin{bmatrix} 2 & -1\\ -1 & 2\end{bmatrix}$ is invertible, since each $\vb{b}$ has exactly
one pair of weights.
\end{example}

\begin{trap}
Dependence is not the same as ``two columns are parallel''. No column of $C$ is a multiple of
another, yet the three columns are dependent because $\vb{u}+\vb{v}+\vb{w}^{*}=\vb{0}$. To
test three vectors you must look for a combination of all of them, not compare them in pairs.
\end{trap}

\subsection{Looking ahead: rectangular matrices}

The habit to build is to look at a matrix and ask what it \emph{does}. $A$ takes differences
and can be undone; $C$ takes cyclic differences, crushes the constant vectors and reaches
only a plane. Matrices need not be square: a $7\times3$ matrix sends $\R^3$ into $\R^7$, can
reach at most a three-dimensional part of $\R^7$, and has no inverse. A matrix that keeps
appearing with rectangular $A$ is $A\T A$: it is square ($3\times3$ here) and symmetric, and
it will be the key to least squares in Part~II (Section~\ref{sec:u17}).

\subsection{Exercises}

\begin{exercise}\label{ex:u02-quiz}
A matrix $A$ has the property that \emph{all} solutions of $A\vb{x}=(1,4,1,1)$ are
\[
  \vb{x}=\begin{bmatrix} 0\\ 1\\ 1\end{bmatrix} + c\begin{bmatrix} 0\\ 2\\ 1\end{bmatrix},
  \qquad c\in\R .
\]
Say everything you can about the columns $\vb{c}_1,\vb{c}_2,\vb{c}_3$ of $A$.
\end{exercise}
\begin{solution}
Since $\vb{x}$ has three components, $A$ has three columns, and since $\vb{b}=(1,4,1,1)$ has
four, each column lies in $\R^4$. Taking $c=0$ gives $A(0,1,1)=\vb{b}$, that is,
$\vb{c}_2+\vb{c}_3=\vb{b}$. Taking $c=1$ gives $A(0,3,2)=\vb{b}$; subtracting,
$A(0,2,1)=\vb{0}$, that is, $2\vb{c}_2+\vb{c}_3=\vb{0}$. Subtract the first relation from the
second: $\vb{c}_2=-\vb{b}$, and then $\vb{c}_3=\vb{b}-\vb{c}_2=2\vb{b}$. So
$\vb{c}_2=(-1,-4,-1,-1)$ and $\vb{c}_3=(2,8,2,2)$.

About $\vb{c}_1$: it cannot be a multiple of $\vb{b}$. If $\vb{c}_1=k\vb{b}$, then
$A(1,k,0)=k\vb{b}-k\vb{b}=\vb{0}$, and $(0,1,1)+(1,k,0)$ would be another solution of
$A\vb{x}=\vb{b}$ whose first component is $1$, which is not of the given form. Conversely,
if $\vb{c}_1$ is not a multiple of $\vb{b}$, then $A\vb{x}=x_1\vb{c}_1+(2x_3-x_2)\vb{b}=\vb{0}$
forces $x_1=0$ and $x_2=2x_3$ (two non-parallel vectors combine to zero only with zero
weights), so the solutions of $A\vb{x}=\vb{0}$ are exactly the multiples of $(0,2,1)$, as
required. Hence: $\vb{c}_1$ is any vector of $\R^4$ that is not a multiple of $\vb{b}$
(in particular not $\vb{0}$), and nothing more can be said.
\end{solution}

\begin{exercise}\label{ex:u02-four}
Let $A_4$ be the $4\times4$ difference matrix, $A_4\vb{x}=(x_1,\,x_2-x_1,\,x_3-x_2,\,x_4-x_3)$.
Write $A_4$ and its inverse $S_4$, and use $S_4$ to solve $A_4\vb{x}=(1,3,5,7)$.
\end{exercise}
\begin{solution}
Solving from the top, $x_1=b_1$, $x_2=b_1+b_2$, $x_3=b_1+b_2+b_3$,
$x_4=b_1+b_2+b_3+b_4$, so
\[
  A_4=\begin{bmatrix} 1&0&0&0\\ -1&1&0&0\\ 0&-1&1&0\\ 0&0&-1&1\end{bmatrix},\qquad
  S_4=A_4^{-1}=\begin{bmatrix} 1&0&0&0\\ 1&1&0&0\\ 1&1&1&0\\ 1&1&1&1\end{bmatrix}.
\]
Then $\vb{x}=S_4(1,3,5,7)=(1,\,4,\,9,\,16)$: running sums of odd numbers are the squares.
Check: $A_4(1,4,9,16)=(1,3,5,7)$.
\end{solution}

\begin{exercise}\label{ex:u02-cyclic}
For the cyclic matrix $C$ of \eqref{eq:u02-C}, decide which of $\vb{b}=(1,0,-1)$,
$(2,-3,1)$ and $(1,1,0)$ lie in the plane of all $C\vb{x}$, and solve $C\vb{x}=\vb{b}$ when it
is possible.
\end{exercise}
\begin{solution}
Use Proposition~\ref{prop:u02-solvable}. For $(1,0,-1)$ the sum is $0$; the particular
solution is $(0,\,b_2,\,b_2+b_3)=(0,0,-1)$, and indeed $C(0,0,-1)=(0+1,\,0-0,\,-1-0)=(1,0,-1)$.
All solutions: $(0,0,-1)+c(1,1,1)$. For $(2,-3,1)$ the sum is $0$; the particular solution
is $(0,-3,-2)$, and $C(0,-3,-2)=(0+2,\,-3-0,\,-2+3)=(2,-3,1)$. All solutions:
$(0,-3,-2)+c(1,1,1)$. For $(1,1,0)$ the sum is $2\neq0$, so this $\vb{b}$ is off the plane
and there is no solution.
\end{solution}

\begin{exercise}\label{ex:u02-which}
Which of these subsets of $\R^3$ are subspaces? (a) all $\vb{b}$ with $b_1+b_2+b_3=0$;
(b) all $\vb{b}$ with $b_1+b_2+b_3=1$; (c) all multiples of $(1,1,1)$; (d) all $\vb{b}$
with $b_1\ge0$; (e) the zero vector alone.
\end{exercise}
\begin{solution}
(a) Yes: it is the plane of all combinations of $\vb{u}$ and $\vb{v}$, and directly, if two
vectors have component sums $0$, so does $c\vb{b}+d\vb{b}'$. (b) No: the zero vector has sum
$0\neq1$, and every subspace contains $\vb{0}$. (c) Yes: a line through the origin;
$c(1,1,1)+d(1,1,1)=(c+d)(1,1,1)$. It is also the set of $\vb{x}$ with $C\vb{x}=\vb{0}$.
(d) No: $(1,0,0)$ is in the set but $-1\cdot(1,0,0)$ is not; the set is closed under
addition but not under multiplication by negative numbers. (e) Yes: $\vb{0}+\vb{0}=\vb{0}$
and $c\,\vb{0}=\vb{0}$; this is the smallest subspace.
\end{solution}
